\section{Aportaciones esperadas}

Del presente trabajo de investigación se proyecta aportar al campo de la Inteligencia Artificial Explicable 
una metodología y una herramienta técnica capaz de cerrar la brecha existente entre los entornos de 
desarrollo de aprendizaje profundo y los ecosistemas gráficos interactivos. Mediante el aprovechamiento de 
las capacidades de Blender y Python, se espera consolidar un flujo de trabajo que no solo facilite la 
visualización e interpretación de datos tensoriales complejos, sino que también optimice su uso en la c
reación de material didáctico e ilustrativo de alta fidelidad espacial.

De manera específica, se contemplan las siguientes aportaciones clave al concluir esta investigación:
\begin{itemize}
    \item \textbf{Un pipeline de automatización e integración:} El código fuente y los scripts de Python 
    necesarios para interceptar y extraer matrices densas de capas convolucionales de producción y 
    traducirlas limpiamente a formatos tridimensionales unificados.
    \item \textbf{Un entorno interactivo de diagnóstico visual (Artifact):} Un sistema procedural desarrollado 
    dentro del editor de nodos de geometría de Blender, diseñado específicamente para realizar operaciones de 
    rebanado vertical, mapas de opacidad y análisis topológico interactivo sobre volúmenes de datos científicos.
    \item \textbf{Impacto en la ingeniería y la investigación:} Una alternativa metodológica viable frente a 
    las herramientas bidimensionales tradicionales que sufren de fragmentación de archivos planos, 
    proveyendo a los desarrolladores de IA de un entorno tridimensional inmersivo para el diagnóstico de 
    filtros y mapas de activación.
    \item \textbf{Recurso para la divulgación científica:} Un marco visual de representación geométrica que
     humanice y desmitifique la naturaleza abstracta del procesamiento interno de la IA, convirtiendo 
     tensores invisibles en material educativo inteligible para un público no especializado.
\end{itemize}